\problemname{Prisoner Challenge}

\begin{quote}
\textit{Note:} in this problem, you have to write your solution in C++.
\end{quote}

In a prison, there are $500$ prisoners. One day, the warden offers them a chance to free themselves.
He places two bags with money, bag A and bag B, in a room. Each bag contains between $1$ and $N$
coins, inclusive. The number of coins in bag A is different from the number of coins in bag B. The
warden presents the prisoners with a challenge. The goal of the prisoners is to identify the bag with
fewer coins.

The room, in addition to the money bags, also contains a whiteboard. A single number must be written
on the whiteboard at any time. Initially, the number on the whiteboard is $0$.

Then, the warden asks the prisoners to enter the room, one by one. The prisoner who enters the room
does not know which or how many other prisoners have entered the room before them. Every time a
prisoner enters the room, they read the number currently written on the whiteboard. After reading the
number, they must choose either bag A or bag B. The prisoner then inspects the chosen bag, thus getting
to know the number of coins inside it. Then, the prisoner must perform either of the following two
actions:
\begin{itemize}
  \item Overwrite the number on the whiteboard with a non-negative integer and leave the room. Note
  that they can either change or keep the current number. The challenge continues after that (unless
  all $500$ prisoners have already entered the room).
  \item Identify one bag as the one that has fewer coins. This immediately ends the challenge.
\end{itemize}

The warden will not ask a prisoner who has left the room to enter the room again.

The prisoners win the challenge if one of them correctly identifies the bag with fewer coins. They lose
if any of them identifies the bag incorrectly, or all $500$ of them have entered the room and not
attempted to identify the bag with fewer coins.

Before the challenge starts, the prisoners gather in the prison hall and decide on a common strategy
for the challenge in three steps.
\begin{itemize}
  \item They pick a non-negative integer $x$, which is the largest number they may ever want to write
  on the whiteboard.
  \item They decide, for any number $i$ written on the whiteboard ($0 \leq i \leq x$), which bag should
  be inspected by a prisoner who reads number $i$ from the whiteboard upon entering the room.
  \item They decide what action a prisoner in the room should perform after getting to know the number
  of coins in the chosen bag. Specifically, for any number $i$ written on the whiteboard
  ($0 \leq i \leq x$) and any number of coins $j$ seen in the inspected bag ($1 \leq j \leq N$), they
  either decide
  \begin{itemize}
    \item what number between $0$ and $x$ (inclusive) should be written on the whiteboard, or
    \item which bag should be identified as the one with fewer coins.
  \end{itemize}
\end{itemize}

Upon winning the challenge, the warden will release the prisoners after serving $x$ more days.

Your task is to devise a strategy for the prisoners that would ensure they win the challenge
(regardless of the number of coins in bag A and bag B). The score of your solution depends on the
value of $x$; see the Scoring section for details.

\section*{Implementation Details}
You should implement the following procedure:

\begin{verbatim}
std::vector<std::vector<int>> devise_strategy(int N)
\end{verbatim}

\begin{itemize}
  \item $N$: the maximum possible number of coins in each bag.
  \item This procedure should return an array $s$ of arrays of $N+1$ integers, representing your
  strategy. The value of $x$ is the length of array $s$ minus one. For each $i$ such that
  $0 \leq i \leq x$, the array $s[i]$ represents what a prisoner should do if they read number $i$
  from the whiteboard upon entering the room:
  \begin{enumerate}
    \item The value of $s[i][0]$ is $0$ if the prisoner should inspect bag A, or $1$ if the prisoner
    should inspect bag B.
    \item Let $j$ be the number of coins seen in the chosen bag. The prisoner should then perform the
    following action:
    \begin{itemize}
      \item If the value of $s[i][j]$ is $-1$, the prisoner should identify bag A as the one with
      fewer coins.
      \item If the value of $s[i][j]$ is $-2$, the prisoner should identify bag B as the one with
      fewer coins.
      \item If the value of $s[i][j]$ is a non-negative number, the prisoner should write that number
      on the whiteboard. Note that $s[i][j]$ must be at most $x$.
    \end{itemize}
  \end{enumerate}
  \item This procedure is called exactly once.
\end{itemize}

Your submission must not contain a \texttt{main} function, and must not read from standard input or
write to standard output: a grader is compiled together with your submission and does that for you.
The declaration above lives in \texttt{prison.h}, which is available both to your submission and as
an attachment, so start your file with \texttt{\#include "prison.h"}.

Any strategy that wins the challenge for every allowed pair of bag contents is accepted. Each test
group puts its own upper bound on $x$, and no group allows $x$ to exceed $500$; see the Scoring
section.

\section*{Constraints}
\begin{itemize}
  \item $2 \leq N \leq 5000$
\end{itemize}

\section*{Scoring}
Your solution will be tested on a set of test groups.
To earn points for a group, you must pass all test cases in that group.

\noindent
\begin{tabular}{| l | l | p{12cm} |}
  \hline
  \textbf{Group} & \textbf{Points} & \textbf{Constraints} \\ \hline
  $1$ & $5$  & $N \leq 500$, and the value of $x$ must not be more than $500$ \\ \hline
  $2$ & $5$  & $N \leq 500$, and the value of $x$ must not be more than $70$ \\ \hline
  $3$ & $90$ & The value of $x$ must not be more than $60$ \\ \hline
\end{tabular}

\medskip
\noindent
If in any of the test cases your strategy does not win the challenge, your score for that group is
$0$.

In group $3$ you can obtain a partial score. Let $m$ be the largest value of $x$ that you print over
all test cases in that group. Your score for group $3$ is calculated according to the following table:

\begin{center}
\begin{tabular}{| l | l |}
  \hline
  \textbf{Condition} & \textbf{Points} \\ \hline
  $40 \leq m \leq 60$    & $20$ \\ \hline
  $26 \leq m \leq 39$    & $25 + 1.5 \times (40 - m)$ \\ \hline
  $m = 25$               & $50$ \\ \hline
  $m = 24$               & $55$ \\ \hline
  $m = 23$               & $62$ \\ \hline
  $m = 22$               & $70$ \\ \hline
  $m = 21$               & $80$ \\ \hline
  $m \leq 20$            & $90$ \\ \hline
\end{tabular}
\end{center}

\section*{Sample Data}
The sample data below is what the judge does: it hands your submission a single line containing $N$,
and the grader it is compiled with prints
\begin{itemize}
  \item line $1$: $x$
  \item line $2 + i$ ($0 \leq i \leq x$): $s[i][0]$ $s[i][1]$ $\ldots$ $s[i][N]$
\end{itemize}
Any strategy that wins the challenge is accepted, so the sample output is only one of many correct
answers.

\section*{Sample Grader}
A sample grader is available under ``attachments'' at the bottom of this page, together with
\texttt{prison.h} and a sample input for it. Unlike the judge, it lets you try your strategy on
scenarios of your own choosing. It reads
\begin{itemize}
  \item line $1$: $N$
  \item line $2 + k$ ($0 \leq k$): $A[k]$ $B[k]$
  \item last line: $-1$
\end{itemize}
Each line except the first and the last one represents a scenario: in scenario $k$ bag A contains
$A[k]$ coins and bag B contains $B[k]$ coins.

The sample grader first calls \texttt{devise\_strategy(N)}. If the returned array does not conform to
the constraints described in Implementation Details, it prints one of the following error messages and
exits:
\begin{itemize}
  \item \texttt{s is an empty array}: $s$ is an empty array (which does not represent a valid
  strategy).
  \item \texttt{s[i] contains incorrect length}: there exists an index $i$ ($0 \leq i \leq x$) such
  that the length of $s[i]$ is not $N+1$.
  \item \texttt{First element of s[i] is non-binary}: there exists an index $i$
  ($0 \leq i \leq x$) such that $s[i][0]$ is neither $0$ nor $1$.
  \item \texttt{s[i][j] contains incorrect value}: there exist indices $i, j$
  ($0 \leq i \leq x$, $1 \leq j \leq N$) such that $s[i][j]$ is not between $-2$ and $x$.
\end{itemize}

Otherwise it produces two outputs. First, it prints
\begin{itemize}
  \item line $1 + k$ ($0 \leq k$): the output of your strategy for scenario $k$ --- the character
  \texttt{A} if applying the strategy leads to a prisoner identifying bag A as the one with fewer
  coins, \texttt{B} if it leads to a prisoner identifying bag B, and \texttt{X} if it does not lead to
  any prisoner identifying a bag.
\end{itemize}
Second, it writes a file \texttt{log.txt} in the current directory in the following format:
\begin{itemize}
  \item line $1 + k$ ($0 \leq k$): $w[k][0]$ $w[k][1]$ $\ldots$
\end{itemize}
The sequence on line $1 + k$ corresponds to scenario $k$ and describes the numbers written on the
whiteboard. Specifically, $w[k][l]$ is the number written by the $(l+1)^{\text{th}}$ prisoner to enter
the room.

\section*{Explanation of Sample}
The sample has $N = 3$. Let $v$ denote the number the prisoner reads from the whiteboard upon entering
the room. One of the correct strategies is as follows:
\begin{itemize}
  \item If $v = 0$ (including the initial number), inspect bag A.
  \begin{itemize}
    \item If it contains $1$ coin, identify bag A as the one with fewer coins.
    \item If it contains $3$ coins, identify bag B as the one with fewer coins.
    \item If it contains $2$ coins, write $1$ on the whiteboard (overwriting $0$).
  \end{itemize}
  \item If $v = 1$, inspect bag B.
  \begin{itemize}
    \item If it contains $1$ coin, identify bag B as the one with fewer coins.
    \item If it contains $3$ coins, identify bag A as the one with fewer coins.
    \item If it contains $2$ coins, write $0$ on the whiteboard (overwriting $1$). Note that this case
    can never happen, as we can conclude that both bags contain $2$ coins, which is not allowed.
  \end{itemize}
\end{itemize}

This is the strategy printed in the sample output, and its value of $x$ is $1$.
